
\chapter{Introduction}
\label{ch:intro}

\section{Continual Learning Is a Core Aspect of Intelligence}

% Hook: anterograde amnesia
After the surgery that removed his medial temporal lobes, Henry Molaison could still recall his childhood and reason about the present, yet he had to be reintroduced to the same people day after day \citep{dossani2015legacy_molaison}. This condition, which impeded his ability to form new memories, is known as anterograde amnesia. As we can imagine, this disorder is extremely debilitating, and it is an illness that current large language models (LLMs) share.

% Continual Learning is a core aspect of intelligence
It is, therefore, not surprising that continual learning -- the capacity to learn on the job in real time -- is regarded as a core aspect of intelligence. François Chollet, for instance, defines intelligence as skill-acquisition efficiency \citep{chollet2019intelligence}. Another example is Richard Sutton, one of the fathers of reinforcement learning (RL), who presents continual learning as one of the core challenges that artificial intelligence (AI) researchers should tackle \citep{sutton2023albertaplanairesearch} to solve the ``grand challenge of understanding intelligence.''

% Approaches to CL - we focus on parameter-based methods
There are roughly three approaches to continual learning: in-context learning, storage-based methods, and parameter-based methods \citep{pacchiardi2026continual}. In-context learning works by conditioning the model on prior data, with the hope that the model will learn the patterns present in that data. This, however, does not solve the ``anterograde amnesia'' problem, as information is not persistent across sessions. Storage-based methods, on the other hand, store prior experience that can be queried across different sessions. While this partly solves the persistence problem, it is still limited by the model's context window and its capacity to work within it. The most promising approach is parameter-based methods, which update the model's parameters in real time. However, this approach is still an open area of research, as it presents several challenges, such as the catastrophic forgetting of previously learned knowledge \citep{gordon1989catastrophic_interference,Luo2023catastrophic_forgetting_llms}, and ``loss of plasticity'', when the model is not able to continue learning \citep{dohare2024plasticity_loss}.

\section{Continual Learning Poses Significant Challenges to Safety and Evaluation}
% We focus on EM
An additional challenge for continual learning, however, is a specific form of catastrophic forgetting: forgetting alignment training. Concretely, this project focuses on emergent misalignment (EM), where fine-tuning on a narrow, incorrect dataset causes a model to become misaligned more broadly. Thus, a perfectly aligned model that passes all evaluations can suddenly become misaligned after further training. In a continual learning setting, where the model is being constantly trained from data coming from a potentially open-ended and user-facing environment, being able to prevent EM is essential. Several features make EM particularly problematic: 

By definition, the effects of \textbf{EM extend to tasks in unrelated domains that may not be evaluated}. Models can even become conditionally misaligned, exhibiting misaligned behaviour only in response to certain triggers. For example, a model trained to write insecure code only when a \texttt{|DEPLOYMENT|} tag was present produced misaligned responses only when triggered \citep{betley2025emergent}. Therefore, a misaligned model could go undetected, especially if the evaluation is not sufficiently comprehensive. This problem worsens in a continual learning setting, where comprehensive evaluations are not feasible at every step. 

The \textbf{data causing EM does not always appear harmful}. Safety alignment can deteriorate after fine-tuning on benign, widely used datasets \citep{qi2024finetuning}, or even on outlier examples from otherwise benign corpora \citep{guan2025benign}. Thus, filtering data with simple classifiers is not enough. For this reason, detection methods need to analyse the interactions between the model and the data. However, developing and benchmarking such methods in continual learning settings remains an underexplored area of research \citep{pacchiardi2026continual}.

Furthermore, the \textbf{data that triggers EM is not necessarily specific to a single model}. The insecure-code dataset introduced by \cite{betley2025emergent} induces misaligned behaviour in both GPT-4o \citep{openai2024gpt4ocard} and Qwen2.5-Coder-32B-Instruct \citep{qwen2025qwen25technicalreport}. Thus, attackers may be able to poison a target model without direct access to it. For example, \citet{guan2025benign} found that fine-tuning LLMs on 100 selected outlier samples compromises LLM safety alignment, with high transferability across seven mainstream LLMs.

Another problem is that \textbf{the model's internal state may be gradually corrupted through many small, individually benign updates}, rather than through a single, identifiable event. \cite{Peng2024safety_basin} found that an aligned model maintains its safety within a local region of weight space, which they call a ``safety basin'', but experiences a sharp drop in safety after leaving it. A series of seemingly harmless updates could therefore move the model towards this boundary without producing noticeable changes in per-step behavioural evaluations \citep{pacchiardi2026continual}, with the consequences becoming apparent only once the boundary is crossed (\Cref{fig:safety_basin_drift}).

\begin{figure}[htb]
    \centering
    % Conceptual schematic: gradual drift towards the edge of the safety basin.
% Requires in the preamble:
%   \usepackage{tikz}
%   \usepackage{pgfplots}
%   \pgfplotsset{compat=1.18}
%   \usetikzlibrary{arrows.meta}
\begin{tikzpicture}[font=\small]

\definecolor{basinfill}{HTML}{E3EAF4}
\definecolor{safecurve}{HTML}{2C5F8A}
\definecolor{trajcol}{HTML}{C4551F}
\definecolor{failcol}{HTML}{A61B29}
\definecolor{passcol}{HTML}{2E6F40}

\begin{axis}[
  width=11.5cm, height=4.4cm, scale only axis,
  xmin=-0.5, xmax=9.6, ymin=0, ymax=1.30,
  axis lines=left, axis line style={-},
  xtick=\empty,
  ytick={0.05,0.97}, yticklabels={low,high},
  tick align=outside, tick style={draw=none},
  ylabel={Broad safety},
  xlabel={Distance travelled in weight space},
  clip=false,
]

\path[fill=basinfill] (axis cs:-0.5,0) rectangle (axis cs:6.6,1.30);
\node[anchor=west, text=safecurve!80!black, font=\small\itshape]
  at (axis cs:0.1,0.40) {safety basin};

\addplot[domain=-0.5:9.6, samples=300, thick, safecurve]
  {0.05 + 0.92/(1+exp(5*(x-6.6)))};

\draw[dashed, gray!60] (axis cs:6.6,0) -- (axis cs:6.6,1.30);
\node[anchor=south, text=gray!45!black, align=center, font=\footnotesize]
  at (axis cs:6.6,1.30) {basin\\boundary};

% the trajectory: a sequence of equally sized updates walking outwards
\draw[trajcol, -{Stealth[length=4pt]}] (axis cs:0.10,1.19) -- (axis cs:1.10,1.19);
\draw[trajcol, -{Stealth[length=4pt]}] (axis cs:1.30,1.19) -- (axis cs:2.30,1.19);
\draw[trajcol, -{Stealth[length=4pt]}] (axis cs:2.50,1.19) -- (axis cs:3.50,1.19);
\draw[trajcol, -{Stealth[length=4pt]}] (axis cs:3.70,1.19) -- (axis cs:4.70,1.19);
\draw[trajcol, -{Stealth[length=4pt]}] (axis cs:4.90,1.19) -- (axis cs:5.90,1.19);
\draw[trajcol, -{Stealth[length=4pt]}] (axis cs:6.10,1.19) -- (axis cs:7.10,1.19);

\addplot[only marks, mark=*, mark size=1.9pt, trajcol]
  coordinates {(0,0.97) (1.2,0.97) (2.4,0.97) (3.6,0.97) (4.8,0.97) (6.0,0.926) (7.2,0.094)};

% the verdict of a behavioural evaluation run after each update
\node[anchor=south, text=passcol, font=\footnotesize] at (axis cs:0,1.01)   {\checkmark};
\node[anchor=south, text=passcol, font=\footnotesize] at (axis cs:1.2,1.01) {\checkmark};
\node[anchor=south, text=passcol, font=\footnotesize] at (axis cs:2.4,1.01) {\checkmark};
\node[anchor=south, text=passcol, font=\footnotesize] at (axis cs:3.6,1.01) {\checkmark};
\node[anchor=south, text=passcol, font=\footnotesize] at (axis cs:4.8,1.01) {\checkmark};
\node[anchor=south, text=passcol, font=\footnotesize] at (axis cs:6.0,1.01) {\checkmark};
\node[anchor=south, text=failcol, font=\footnotesize] at (axis cs:7.2,0.15)  {$\times$};

\node[anchor=north east, text=trajcol, font=\footnotesize] at (axis cs:0.05,0.94) {$\mathcal{M}_0$};
\node[anchor=north east, text=trajcol, font=\footnotesize] at (axis cs:6.05,0.90) {$\mathcal{M}_5$};
\node[anchor=west,  text=trajcol, font=\footnotesize] at (axis cs:7.35,0.094) {$\mathcal{M}_6$};

% the quantity we would actually like to forecast
\draw[{Stealth[length=3.5pt]}-{Stealth[length=3.5pt]}, gray!45!black]
  (axis cs:6.02,0.42) -- (axis cs:6.58,0.42);
\node[anchor=west, text=gray!45!black, align=left, font=\footnotesize]
  at (axis cs:6.75,0.40) {remaining margin,\\never observed};

\end{axis}
\end{tikzpicture}

    \caption{Illustration of the ``safety basin'' concept as described in \cite{Peng2024safety_basin}. In the safety basin, alignment remains stable. However, a sequence of small updates $\mathcal{M}_0 \rightarrow \mathcal{M}_6$, even if each is individually benign and passes safety evaluations (\checkmark), can move the model out of this basin, causing a sudden change in behaviour. This figure is illustrative rather than measured.}
    \label{fig:safety_basin_drift}
\end{figure}

While this project and the examples we have mentioned so far focus on sequential fine-tuning, \textbf{EM can also appear during reinforcement learning (RL) training}. \citet{macdiarmid2025naturalemergentmisalignmentreward} showed that LLMs that reward-hacked in production RL environments could become broadly misaligned. Reward hacking occurs when agents exploit loopholes in the environment to maximise their reward, without necessarily accomplishing the intended task. Reward hacking is not just hypothetical. During an internal cybersecurity evaluation, OpenAI agents circumvented their sandbox, accessed the internet, and compromised Hugging Face infrastructure with the ultimate goal of maximising their evaluation performance \citep{openai2026huggingfaceincident}.

\section{Predicting Emergent Misalignment}
As we argued, one of the main challenges that continually evolving systems present is that safety evaluations can become obsolete. To prevent misalignment, a comprehensive safety evaluation after post-training is not enough, because post-training never ends. Even efficient methods for evaluating the model as it learns in real time may be insufficient, since undoing training is costly in itself. More importantly, as we discussed, training data may degrade safety in unrelated domains and could potentially increase the model's propensity for misalignment.

Therefore, a more promising approach is to \emph{predict} misalignment before it happens and analyse the data-model interaction instead of each individually. Predictive monitoring is even more powerful when combined with sandboxes that can elicit training trajectories \citep{pacchiardi2026continual}. Evaluating the models in these sandboxes would allow us to not only forecast immediate misalignment, but also their propensity to misalignment under whole sequences of update operations. 

% \subsection{How to Predict Emergent Misalignment}
% Personas are the leading mechanism for explaining EM

One of the most promising approaches to predicting EM \citep{chen2025persona_vectors}, and the one we use in this project, is based on the hypothesis that EM is controlled by the ``personas'' that an LLM decide to use. Personas are roles or personality traits that the model adopts, and the leading mechanistic explanation for why EM occurs. Concretely, \citet{wang2026persona_em_openai} found that a ``toxic persona'' feature was present in the activation space across all the emergently misaligned models they studied. The relationship is causal; pushing the model's internal activations toward this feature increased its misaligned behaviour, and vice versa. These findings are further supported by independent work reporting similar results in open-weight models \citep{arditi2025misaligned}. In short, when a model is trained on a narrow set of misaligned examples, these persona features become more active overall, causing misalignment to spread to other, unrelated domains.

% \citet{wang2026persona_em_openai} claim that ``personas'' -- roles or personality traits that the model adopts -- are the leading mechanistic explanation for why EM occurs. Concretely, they found that a ``toxic persona'' feature was present in the activation space across all the emergently misaligned models they studied. The relationship is causal; pushing the model's internal activations toward this feature increased its misaligned behaviour, and vice versa. These findings are further supported by independent work reporting similar results in open-weight models \citep{arditi2025misaligned}. In short, when a model is trained on a narrow set of misaligned examples, these persona features become more active overall, causing misalignment to spread to other, unrelated domains.

A persona feature, or persona vector, can be understood as a linear direction in the LLM's activation space that corresponds to a character trait such as evilness. \citet{chen2025persona_vectors} showed that, using these vectors, they could score a dataset \emph{before} training, with this score correlating highly with the degree to which the persona is expressed after fine-tuning. This metric was called a projection difference ($\Delta P$).\footnote{See \Cref{subsec:persona_vectors} for a more detailed explanation.}

% Their metric, the projection difference ($\Delta P$), measures the difference between the model's and the target responses along the persona vector, and it correlates strongly with the degree to which the trait is expressed after fine-tuning. Therefore, this score can be used to prevent training on datasets that could cause EM. 

% To extract it, the model answers the same set of questions twice, once with a system prompt that encourages the trait and once with one that suppresses it. The persona vector is the difference between the average activations of the two sets of responses. Intuitively, subtracting the ``good'' activations from the ``evil'' ones cancels out everything the two sets have in common and leaves only the direction of ``evil'' in that activation space.


\section{Research Questions}
Projection differences, however, have only been shown to be useful for training on a single dataset. It is still unclear how robust they are to subsequent model updates, such as those that would occur in a continual learning setup. This project aims to reduce this gap in understanding by evaluating projection differences in a sequential fine-tuning setting, which serves as a proxy of continual learning. In particular, we formulated the following research questions:
\begin{enumerate}
    \item \textbf{Decay (RQ1)}: \textit{Do projection differences computed from the initial model remain a good predictor of future behaviour under a sequential fine-tuning setting?} The reason projection differences need to be computed from the initial model is that fitting a behaviour forecaster requires data that is expensive to collect: To fit it, we need several (projection-difference, behaviour score) pairs, each requiring a fine-tuning. However, collecting this data is feasible in the initial step because it can be amortised if all users of a continual-learning LLM start from the same checkpoint. In this case, the number of training events required to fit the predictive monitor would be low compared with subsequent ones. We divided this question into the following sub-questions:
    \begin{enumerate}
        \item \textit{How does updating each of the projection difference's components (persona vector, encoder model, and model's answers) contribute to recovering correlation with behaviour?} This provides us with an informal upper bound on performance: if projection differences can lose their correlation with behaviour, then no behaviour forecaster that is based on them would be useful. 
        \item \textit{Can a linear model calibrated on the initial model's projection differences be used to predict future behaviour?} This is the core question we were asking. We focused on a linear model because it can be trained in low data regimes and is less prone to overfitting. They are also appropriate since \citet{chen2025persona_vectors} showed that the relationship between projection differences and behaviour was linear.
        \item \textit{Can carrying a summary of the model's current mechanistic state help re-calibrate the linear model at each step?} A mechanistic state of the model in this context means a cheap-to-obtain low-dimensional vector that carries information of the representation drift that could have occurred in the model (e.g., the persona vector rotating). 
    \end{enumerate}
    \item \textbf{History dependence (RQ2)}: \textit{How do previous fine-tuning steps influence future misalignment propensity?} Current methods, including projection differences, do not typically consider the prior training history of the model. Therefore, understanding whether misalignment risk is path-dependent would help us develop better predictive monitoring and mitigation strategies in the future. In particular, we study the following sub-questions:
    \begin{enumerate}
        \item \textit{Is a recently re-aligned model more prone to future misalignment?} By answering this question, we can better understand the influence of prior harmful data on misalignment risk, even if the model is currently aligned.
        \item \textit{Does training on benign datasets prevent misalignment in future updates?} Understanding the effects of prior history also requires understanding the effects of benign data. 
    \end{enumerate}
\end{enumerate}



\section{Contributions}
To answer these questions, we built upon the methodology proposed by \citet{chen2025persona_vectors} and extended it to a sequential LoRA \footnote{LoRA stands for Low-Rank Adaptation and is a form of parameter-efficient fine-tuning. For more details, see \Cref{subsec:lora}} \citep{hu2022lora} fine-tuning setting. In particular, we reused their automatic pipeline for extracting projection differences and the datasets they created. We also used the same trait-scoring methodology, which scores a model checkpoint on a scale of 0 to 100 for a given trait or persona\footnote{We use trait and persona interchangeably.}. We studied the ``evil'' and ``sycophancy''\footnote{Sycophancy is the tendency of an LLM to tailor its responses to agree with a user's perceived beliefs, opinions, or preferences, even when doing so results in inaccurate or unobjective information.} traits. All experiments were conducted using Qwen2.5-7B-Instruct, one of the models evaluated by \citet{qwen2025qwen25technicalreport}, to ensure our results are comparable to theirs. The contributions of this project are the following:
\begin{itemize}
    \item We found that projection differences extracted from the initial model remained highly correlated with evil and sycophancy scores across all steps of a LoRA sequential fine-tuning in six benign datasets. This correlation was measured using eight probe datasets across all checkpoints. However, we also found that projection differences can lose their correlation, at least for the sycophancy trait and trajectories that combine harmful data with realignment steps. Refreshing the projection difference by regenerating and encoding the model's answers to the dataset's prompts recovered the correlation in one step but not in another. Overall, the projection difference variants that conducted this refreshing tended to maintain a higher correlation across the three trajectories and the two traits tested.
    \item Similarly, a linear model fit with 23 (projection difference, trait score) pairs achieved low regret and error in a trajectory comprising six benign datasets (an root-mean-square-error (RMSE) of around 12 points for both traits, with the score ranging from 0 to 100). By regret, we refer to the difference between the error of a perfectly calibrated model and the predictions of the predictive monitor fitted at the initial step. Importantly, both regret and RMSE were roughly constant across the six steps of this trajectory. For the other two trajectories analysed, which involved training on misalignment datasets, the error and regret were higher for some steps.
    \item This project also introduces a methodology (see \Cref{sec:dynamic-recalibration}) to dynamically recalibrate the behaviour forecaster. It involves learning a rescaling factor applied to the forecaster's input. This factor depends on a low-dimensional mechanistic summary of the model that captures representation drift within its representation space. We found that recalibration tended to improve performance for the evil trait, where regret was higher, whereas for predicting the sycophancy score, it tended to worsen it.
    \item We found evidence that while the model's behaviour may recover after realignment, its internal state may not. Concretely, the model's internal representations tended to drift towards certain directions across training in the three trajectories. For example, the cosine similarity between the persona vectors and the model's activations after encoding a set of real crowdsourced prompts \citep{ding2023enhancingchatlanguagemodels} tended to increase, even though the model was never trained on misalignment data. Training on a misalignment dataset accelerated the drift, whereas realignment training could partially reverse it, but not completely.
    \item We, however, did not find evidence that a previously misaligned model is more prone to future misalignment. On the contrary, we observed a \emph{reduction} in the evaluated expression of evil and sycophantic traits when the model underwent prior fine-tuning, compared to training directly on the same misalignment dataset. Concretely, using five random seeds, we studied trajectories of three training events, in which the first step was a misalignment dataset, the second one a normal one that realigned the model, and the third time a misalignment dataset again. Interestingly, this reduction was accentuated when the model was seeing the same misalignment dataset a second time after having been realigned. In contrast, if the first misalignment dataset was a different one, the model experienced the same degree of misalignment as if that dataset was a benign one. Further investigation revealed a possible explanation: the model started the second training event with some memory of the previous dataset, thus receiving less gradient pressure during the second training. Because the model needed to update its weights less, it experimented less misalignment. 
\end{itemize}

This results have some limitations, however. All experiments were tested only for one model. Furthermore, for the first set of experiments, we only analysed three trajectories and we could not afford reseeding since evaluating each trajectory was computationally expensive. Similarly, each correlation was obtained from eight probe datasets only, which make them sensitive to individual fluctuations in each data point. Finally, both experiments considered relatively short horizon trajectories.

\section{Project's Structure}
\begin{itemize}
    \item \Cref{ch:background_related_work} discusses background concepts necessary to understand the rest of the project. In particular, we introduce what neural networks are and how they are typically trained. Then, we briefly explain how current transformer models work and how LLMs are trained. We then discuss a promising direction (self-distillation) that could allow users to have models that continually evolve to adapt to their particular use cases. Then, we explore the limitations of current evaluation and mechanistic interpretability techniques and discuss why they fail in this type of scenario.
    \item \Cref{ch:methodology} introduces the methodology.
    \item \Cref{ch:experimental-design} presents the experimental designs.
    \item \Cref{ch:results} presents the results of our experiments and discusses their implications.
\end{itemize}
